%% ===========================================================================
\FloatBarrier
\section{Parameter identification and limitations}
\label{sec:calibration}
%% ===========================================================================

Four parameters are supplied by no constitutive table: the window size $a$,
the fractional order $\alpha$, the relaxation time $\tau$ and the flexoelectric
coefficient $\mu$. This section states how each is obtained from a separate
observable and over which ranges the results of Sec.~\ref{sec:results} hold, so
that a reader can tell a measured input from an assumed one.

\FloatBarrier
\subsection{Measurement routes for the four parameters}
\label{sec:calib_a_mu}

The four are not of the same kind. The geometry of the interaction, $M_{mn}$
and the orientation $\theta$ of its principal axes, comes from size-effect and
texture measurements; the material response, $\alpha$ and $\tau$, from
dielectric loss; $\mu$ from a bending or a direct charge measurement. The
protocol is ordered so that each step fixes one quantity and the later steps
verify the parameters fixed in preceding steps, as in the
characteristic-length identification of the one-scale theory~\cite{fem1}.

\begin{enumerate}
\item Bending of strips cut along $x$ and $y$ at three or more thicknesses in
  $h/\sqrt{M_{22}}\in[2,20]$: the ratio of apparent to classical compliance
  against $h$ fixes $M_{22}$ through the thickness.
\item The same strips cut at $0^\circ,15^\circ,\dots,90^\circ$: the compliance
  and the first frequency against the cut angle fix $M_{11}/M_{22}$ and the
  orientation $\theta$.
\item Ultrasonic dispersion: the phase velocity $c(k)$ over the band in which
  the gradient term is visible fixes $\tfrac14 M_{mn}$ without any boundary
  condition.
\item A boundary-layer profile by digital image correlation: the decay length
  of the strain at a clamped edge fixes $M_{22}$ from one specimen.
\item An image of the microstructure: integrating $s_ms_n$ over the
  reconstructed window fixes the anisotropy of $M$ directly.
\item Specimens spanning the aspect ratio and the window angle: the coupling
  factor and the first frequency against both are then a prediction, not a fit.
\item The reconstructed two-point correlation: its decay lengths set the
  adequacy of the second-moment description, which this paper assumes and does
  not test.
\end{enumerate}

Steps 1, 3 and 4 determine $M_{22}$ from three independent observables, so
agreement among them checks the second-moment description instead of repeating
one measurement; step~6 turns the calibrated tensor into a prediction of
Figs.~\ref{fig:orientation} and~\ref{fig:design}; step~7 is the test of the
description itself and is the open problem of the model.

\paragraph{The flexoelectric coefficient} The direct route is a four-point
bending test on a poled bar with electrodes on the two faces: for the geometry
of Fig.~\ref{fig:geometry}(a) the open-circuit voltage at a known curvature
$\kappa_{c}$ is
\begin{equation}
  V = \frac{\mu\,\kappa_{c}\,h}{\kappa_{33}},
  \label{eq:muV}
\end{equation}
with no discretization involved and $\kappa_{33}=4.9583\times10^{-10}$\,\si{
\farad\per\meter} the polar-axis permittivity of \eqref{eq:D}. The inverse route
measures the charge released by a device whose curvature is known from its
deflection. Both fix the magnitude, which is linear in $\mu$; the sign comes
from the crystal symmetry. The value used here,
$\mu=10^{-6}$\,\si{\coulomb\per\meter}, is the order reported for perovskite
ceramics, and Sec.~\ref{sec:material} states once how every flexoelectric
coupling below follows it.

\FloatBarrier
\subsection{Identification of the fractional order and of the relaxation time}
\label{sec:calib_alpha_tau}

The pair $(\alpha,\tau)$ comes from a dielectric loss measurement. The
fractional Kelvin--Voigt element \eqref{eq:kvfreq} has
\begin{equation}
  \tan\delta=\frac{(\omega\tau)^{\alpha}\sin(\pi\alpha/2)}
     {1+(\omega\tau)^{\alpha}\cos(\pi\alpha/2)},
  \label{eq:tandelta}
\end{equation}
which increases monotonically to $\tan(\pi\alpha/2)$ and has no maximum, and
both $|K^{*}|/K$ and $\tan\delta$ are monotone in $\omega\tau$; so no single
frequency carries $\tau$, and the pair $(|K^{*}|/K,\tan\delta)$ over two
decades is fitted with $(\alpha,\tau)$. At $\omega\tau=1$ the loss is
$\tan(\pi\alpha/4)$, which gives $\alpha$ from one decade once $\tau$ is known.
The ring-down gives the same information in the time domain, but not from one
number: for small loss
\begin{equation}
  \zeta=\frac{(\omega\tau)^{\alpha}\sin(\pi\alpha/2)}{2}
  \quad\Longrightarrow\quad
  \delta\approx\pi(\omega\tau)^{\alpha}\sin(\pi\alpha/2),
  \label{eq:zetasmall}
\end{equation}
so at $\alpha=1$, $\delta=\pi\omega\tau$, which reproduces \eqref{eq:zeta} to
the $1.2\,\%$ of the small-loss approximation. A single decrement fixes only
that product, so the pair needs two frequencies or one decrement at a known
$\omega\tau$.

Two cautions follow from Sec.~\ref{sec:memorystudy}: the calibration must say
which matrix the memory multiplies, since one measured decrement gives two
different $(\alpha,\tau)$ pairs under $K_{\rm eff}$ and $K_{uu}$ (their damping
ratios differ by the factor $0.450$ at $\alpha=1$); and the fit is valid only
where the measurement is -- a broadened single-power-law loss. Two well
separated loss peaks need two elements.

\FloatBarrier
\subsection{Positioning}
\label{sec:positioning}

Three lines of work meet in this paper and none of them carries all of its
ingredients. The one-scale gradient theory~\cite{fem1} supplies the $C^1$
Bogner--Fox--Schmit technology and the geometric meaning of the length tensor,
for an elastic solid without electromechanical coupling. The
polarization-gradient theory of Mindlin~\cite{Mindlin1968} and the
flexoelectric implementations that followed it~\cite{ChenYaoYvonnet2023,
JoshanSantapuri2024,Wang2024PhaseField} supply the electromechanical coupling,
with a scalar or a fitted gradient length and without dielectric memory.
Fractional constitutive models~\cite{BagleyTorvik1983,Koeller1984,
EnelundLesieutre1999} supply the broad relaxation spectrum, in mechanical
settings that have no strain gradient and no electrodes. What this paper adds
is the combination and the two consequences of the combination: the window
enters as a tensor that can be measured from a microstructure, the memory
matrix becomes an explicit modeling choice with a measurable consequence
($0.450$ on the damping ratio at $\alpha=1$), and the coupling factor of a
strain-gradient electromechanical strip is separated into a material property
($0.5581$, thickness independent) and a size effect (a factor $24$).

\FloatBarrier
\subsection{Validity window}
\label{sec:limit_window}

Three inequalities bound the results, and each is a modeling assumption rather
than a numerical tolerance.

\begin{enumerate}
\item \emph{Small coupling at the operator level.} \eqref{eq:kstatic} is used
  with the uncoupled denominator only where $k^{2}\ll1$; here
  $k^{2}=0.5581$, which is \emph{not} small, so the coupled denominator is used
  and the uncoupled form appears only in the comparison with the closed form.
\item \emph{Small window.} The truncation \eqref{eq:Mdef} is written for a
  window small against the length over which the strain gradient varies, and
  for the strip of Table~\ref{tab:params} we take $a/L\le0.2$. That is an
  assumed operating range, not a validated truncation threshold: with no
  experiment and no exact averaging operator available, nothing here fixes
  where the truncation fails. What Fig.~\ref{fig:window}(a) shows is that
  inside the range the retained second moment reproduces the exact average to
  a fraction of a per cent on every mesh, which makes the model
  \emph{self-consistent} there and is all that is claimed.
\item \emph{Operating band.} The damping map is claimed for $\omega\tau\le1$,
  the band in which the spring-pot relaxes around the device frequency. Beyond
  it the map extrapolates the fitted law, and the integer-order alternative
  departs from it by the factors quoted in Sec.~\ref{sec:memorystudy}.
\end{enumerate}

\FloatBarrier
\subsection{Limitations}
\label{sec:limitations}

The formulation is two-dimensional with a plane-stress reduction, which is a
modeling choice: plane strain would change the extensional coupling to
$0.5132$, and the three-dimensional solution differs from both, so the
thin-strip results of Sec.~\ref{sec:thickness} are two-dimensional results. Of
the inputs, two are not material constants of this crystal. The density
rescales every frequency by a common factor, so the kilohertz numbers of
Secs.~\ref{sec:orientation} and~\ref{sec:application} move while the ratios
stand. The flexoelectric coefficient scales every flexoelectric coupling as
$\mu^{2}$, so the growth factor, the saturation and the comparison with the
piezoelectric response are unaffected while the absolute coupling is not.

The precision of Sec.~\ref{sec:thickness} is a limitation as much as a result:
two significant digits of that coupling are resolved, and the same holds, less
visibly, for any difference of nearly equal eigenvalues. The verification of
Secs.~\ref{sec:verification} and~\ref{app:verification} shows that the
implementation solves the stated equations; it does not show that they
describe a particular material, which needs the measurements of
Secs.~\ref{sec:calib_a_mu} and~\ref{sec:calib_alpha_tau}.

\FloatBarrier
